\section{Banco de questões — Língua Portuguesa}

\subsection*{N1 — Fundamento competitivo}

\begin{questao}{\nivel{N1} 01.01 — Cebraspe/Certo ou Errado}
No trecho ``Embora o sistema tenha reduzido o tempo médio de atendimento, aumentou o número de reaberturas'', a substituição de \emph{Embora} por \emph{Porque} preservaria a correção gramatical, mas alteraria a relação lógica entre as orações.
\end{questao}

\begin{questao}{\nivel{N1} 01.02 — Cebraspe/Certo ou Errado}
A frase ``Nem todos os processos foram revisados'' é semanticamente equivalente a ``Nenhum processo foi revisado''.
\end{questao}

\begin{questao}{\nivel{N1} 01.03 — Cebraspe/Certo ou Errado}
Em ``O gestor informou ao auditor que seu relatório estava incompleto'', o pronome possessivo pode gerar ambiguidade referencial.
\end{questao}

\begin{questao}{\nivel{N1} 01.04 — Cebraspe/Certo ou Errado}
Em ``Há inconsistências nos registros'', a substituição de \emph{há} por \emph{existem} preserva o sentido básico e exige flexão no plural do novo verbo.
\end{questao}

\begin{questao}{\nivel{N1} 01.05 — Cebraspe/Certo ou Errado}
Em ``Devem haver soluções alternativas'', a flexão plural está de acordo com a norma-padrão, porque o termo \emph{soluções} funciona como sujeito de \emph{haver}.
\end{questao}

\begin{questao}{\nivel{N1} 01.06 — Cebraspe/Certo ou Errado}
A inserção de vírgulas em ``Os servidores que concluíram o curso receberam certificado'' para formar ``Os servidores, que concluíram o curso, receberam certificado'' pode alterar o alcance da informação.
\end{questao}

\begin{questao}{\nivel{N1} 01.07 — Cebraspe/Certo ou Errado}
Em ``referiu-se à norma'', o sinal indicativo de crase decorre da combinação da preposição exigida pelo verbo com o artigo feminino que antecede \emph{norma}.
\end{questao}

\begin{questao}{\nivel{N1} 01.08 — Cebraspe/Certo ou Errado}
Em ``não se verificaram falhas'', a presença da palavra negativa favorece a próclise do pronome átono.
\end{questao}

\begin{questao}{\nivel{N1} 01.09 — Cebraspe/Certo ou Errado}
A substituição de \emph{retificar} por \emph{ratificar} preserva o sentido em ``o setor deverá retificar os dados incorretos''.
\end{questao}

\begin{questao}{\nivel{N1} 01.10 — Cebraspe/Certo ou Errado}
Na frase ``Se a equipe viesse hoje, concluiríamos a análise'', há correlação adequada entre o imperfeito do subjuntivo e o futuro do pretérito.
\end{questao}

\subsection*{N2 — Aplicação}

\begin{questao}{\nivel{N2} 01.11 — FGV/interpretação}
Considere: ``A implantação reduziu o tempo de atendimento, mas elevou a taxa de reabertura.'' É inferência autorizada pelo trecho:
\begin{enumerate}[label=\Alph*)]
\item a implantação fracassou integralmente;
\item a implantação melhorou uma métrica e piorou outra;
\item a reabertura cresceu porque o atendimento ficou mais rápido;
\item os usuários ficaram mais insatisfeitos;
\item a implantação deve ser descontinuada.
\end{enumerate}
\end{questao}

\begin{questao}{\nivel{N2} 01.12 — FCC/coesão}
Em ``O relatório foi devolvido ao setor, pois continha inconsistências'', a conjunção \emph{pois} introduz, no contexto:
\begin{enumerate}[label=\Alph*)]
\item consequência;
\item causa/explicação;
\item condição;
\item concessão;
\item finalidade.
\end{enumerate}
\end{questao}

\begin{questao}{\nivel{N2} 01.13 — Cebraspe/Certo ou Errado}
No período ``Caso os dados estejam completos, o relatório será emitido'', a substituição de \emph{Caso} por \emph{Embora} preservaria a estrutura gramatical geral, mas mudaria condição para concessão.
\end{questao}

\begin{questao}{\nivel{N2} 01.14 — FGV/pontuação}
Assinale a frase em que a vírgula está empregada de acordo com a norma-padrão.
\begin{enumerate}[label=\Alph*)]
\item Os auditores, verificaram os documentos.
\item A equipe analisou, os registros do contrato.
\item Ao final da inspeção, a equipe registrou as evidências.
\item O relatório concluiu que, havia falhas graves.
\item Foram identificados, riscos e controles deficientes.
\end{enumerate}
\end{questao}

\begin{questao}{\nivel{N2} 01.15 — FCC/concordância}
Assinale a opção correta segundo a norma-padrão.
\begin{enumerate}[label=\Alph*)]
\item Fazem dois anos que o sistema foi implantado.
\item Devem existir alternativas de solução.
\item Devem haver documentos arquivados.
\item Houveram falhas relevantes.
\item Existe inconsistências no cadastro.
\end{enumerate}
\end{questao}

\begin{questao}{\nivel{N2} 01.16 — FGV/regência}
Assinale a frase correta.
\begin{enumerate}[label=\Alph*)]
\item O auditor assistiu o julgamento, no sentido de presenciá-lo.
\item O servidor obedeceu o regulamento.
\item A equipe preferiu revisar os dados do que arquivá-los.
\item O gestor informou o resultado ao servidor.
\item O candidato aspirava o cargo, no sentido de almejá-lo.
\end{enumerate}
\end{questao}

\begin{questao}{\nivel{N2} 01.17 — Cebraspe/Certo ou Errado}
Em ``O gestor informou o resultado ao servidor'', a substituição por ``O gestor informou o servidor do resultado'' preserva estrutura de regência admitida na norma-padrão.
\end{questao}

\begin{questao}{\nivel{N2} 01.18 — FCC/crase}
Assinale a opção em que o emprego da crase é obrigatório.
\begin{enumerate}[label=\Alph*)]
\item Entregou o documento a Vossa Excelência.
\item Começou a revisar os autos.
\item Referiu-se a uma norma interna.
\item Dirigiu-se à unidade responsável.
\item Ficou frente a frente com o gestor.
\end{enumerate}
\end{questao}

\begin{questao}{\nivel{N2} 01.19 — FGV/colocação pronominal}
Assinale a opção mais adequada à norma-padrão formal.
\begin{enumerate}[label=\Alph*)]
\item Não verificou-se o erro.
\item Não se verificou o erro.
\item Se verificou o erro ontem, em início absoluto de oração formal.
\item Jamais constatou-se a falha.
\item Quem informou-me o resultado?
\end{enumerate}
\end{questao}

\begin{questao}{\nivel{N2} 01.20 — Cebraspe/Certo ou Errado}
Em ``O relatório que foi publicado contém ressalvas'', o termo \emph{que} é pronome relativo e retoma \emph{relatório}.
\end{questao}

\begin{questao}{\nivel{N2} 01.21 — FCC/sintaxe}
Em ``O auditor afirmou que os dados estavam completos'', a oração introduzida por \emph{que} exerce função equivalente a:
\begin{enumerate}[label=\Alph*)]
\item adjunto adverbial de causa;
\item complemento do verbo \emph{afirmou};
\item oração adjetiva restritiva;
\item aposto explicativo;
\item vocativo.
\end{enumerate}
\end{questao}

\begin{questao}{\nivel{N2} 01.22 — FGV/reescrita}
A frase ``A comissão analisou o processo e emitiu parecer'' mantém o sentido básico em:
\begin{enumerate}[label=\Alph*)]
\item O parecer foi emitido pela comissão após a análise do processo.
\item A comissão emitiu parecer, embora não tenha analisado o processo.
\item Como emitiu parecer, a comissão não analisou o processo.
\item A análise do processo impediu a emissão de parecer.
\item O processo analisou a comissão antes do parecer.
\end{enumerate}
\end{questao}

\begin{questao}{\nivel{N2} 01.23 — Cebraspe/Certo ou Errado}
A passagem da voz ativa ``A equipe revisou os dados'' para a passiva ``Os dados foram revisados pela equipe'' preserva os participantes básicos do evento, embora altere o foco informacional.
\end{questao}

\begin{questao}{\nivel{N2} 01.24 — FCC/ortografia}
Assinale a opção corretamente grafada no contexto.
\begin{enumerate}[label=\Alph*)]
\item O setor irá ratificar os erros de digitação.
\item O risco é eminente e ocorrerá em poucos minutos.
\item Não se sabe por que os dados divergiram.
\item O documento foi entregue há dois dias atrás.
\item O servidor agiu mau durante a reunião.
\end{enumerate}
\end{questao}

\begin{questao}{\nivel{N2} 01.25 — Cebraspe/Certo ou Errado}
Na frase ``Há dois anos o órgão implantou o sistema'', a substituição de \emph{há} por \emph{a} produziria erro na expressão de tempo decorrido.
\end{questao}

\subsection*{N3 — Integração de concurso}

\begin{questao}{\nivel{N3} 01.26 — FGV/reescrita e sentido}
Original: ``Apenas os contratos de alto risco serão examinados inicialmente.'' A reescrita que preserva o sentido é:
\begin{enumerate}[label=\Alph*)]
\item Todos os contratos serão examinados, sobretudo os de alto risco.
\item Inicialmente, serão examinados somente os contratos de alto risco.
\item Nenhum contrato de baixo risco será examinado em qualquer momento.
\item Os contratos de alto risco serão os únicos examinados definitivamente.
\item Alguns contratos de alto risco não serão examinados inicialmente.
\end{enumerate}
\end{questao}

\begin{questao}{\nivel{N3} 01.27 — Cebraspe/Certo ou Errado}
A substituição de ``ao menos um contrato não foi revisado'' por ``nenhum contrato foi revisado'' altera o quantificador e torna a segunda afirmação logicamente mais forte.
\end{questao}

\begin{questao}{\nivel{N3} 01.28 — FCC/pontuação e sentido}
Compare:

I. Os gestores que não apresentaram justificativa foram notificados.

II. Os gestores, que não apresentaram justificativa, foram notificados.

A diferença é que:
\begin{enumerate}[label=\Alph*)]
\item I e II são semanticamente idênticas;
\item I restringe o grupo notificado aos gestores sem justificativa; II apresenta a falta de justificativa como explicação atribuída ao conjunto referido;
\item II restringe e I explica;
\item as vírgulas são facultativas e sem efeito;
\item II é obrigatoriamente agramatical.
\end{enumerate}
\end{questao}

\begin{questao}{\nivel{N3} 01.29 — FGV/conectores}
Original: ``O sistema é estável, embora apresente limitações de escala.'' Qual reescrita preserva a relação concessiva?
\begin{enumerate}[label=\Alph*)]
\item O sistema é estável porque apresenta limitações de escala.
\item Apesar de apresentar limitações de escala, o sistema é estável.
\item O sistema é estável, portanto apresenta limitações.
\item Se apresentar limitações, o sistema será estável.
\item Como é estável, o sistema apresenta limitações.
\end{enumerate}
\end{questao}

\begin{questao}{\nivel{N3} 01.30 — Cebraspe/Certo ou Errado}
Em ``Ainda que os dados estejam incompletos, a análise continuará'', o modo subjuntivo é compatível com a relação concessiva introduzida por \emph{ainda que}.
\end{questao}

\begin{questao}{\nivel{N3} 01.31 — FCC/regência e crase}
Assinale a frase correta.
\begin{enumerate}[label=\Alph*)]
\item A equipe aspirava à melhoria do processo.
\item A equipe aspirava à um resultado melhor.
\item O auditor obedeceu à regulamento interno.
\item O servidor preferiu a revisão do que o arquivamento.
\item O gestor assistiu à o julgamento.
\end{enumerate}
\end{questao}

\begin{questao}{\nivel{N3} 01.32 — FGV/ambiguidade}
A melhor reescrita para eliminar a ambiguidade de ``O chefe informou ao analista que seu parecer seria revisto'', caso o parecer pertença ao analista, é:
\begin{enumerate}[label=\Alph*)]
\item O chefe informou ao analista que seu parecer seria revisto.
\item O chefe informou que o parecer seria revisto ao analista.
\item O chefe informou ao analista que o parecer deste seria revisto.
\item O chefe informou ao analista sobre seu parecer.
\item Seu parecer seria revisto pelo chefe e pelo analista.
\end{enumerate}
\end{questao}

\begin{questao}{\nivel{N3} 01.33 — Cebraspe/Certo ou Errado}
Na norma-padrão, ``Fazem cinco meses que a auditoria começou'' deve ser corrigido para ``Faz cinco meses...'', porque \emph{fazer} indicando tempo decorrido é impessoal.
\end{questao}

\begin{questao}{\nivel{N3} 01.34 — FCC/concordância}
Assinale a opção correta.
\begin{enumerate}[label=\Alph*)]
\item É necessário a revisão dos autos.
\item É necessária a revisão dos autos.
\item São necessário revisar os autos.
\item É necessários os documentos.
\item Fazem necessária a revisão.
\end{enumerate}
\end{questao}

\begin{questao}{\nivel{N3} 01.35 — FGV/coesão}
No trecho ``O relatório apontou falhas. Essas inconsistências, contudo, não impediram a continuidade do serviço'', os termos \emph{Essas inconsistências} e \emph{contudo} exercem, respectivamente:
\begin{enumerate}[label=\Alph*)]
\item referenciação anafórica e relação adversativa;
\item catáfora e causa;
\item elipse e condição;
\item repetição sem coesão e conclusão;
\item nominalização e finalidade.
\end{enumerate}
\end{questao}

\begin{questao}{\nivel{N3} 01.36 — Cebraspe/Certo ou Errado}
Em ``O relatório foi publicado; portanto, tornou-se acessível ao público'', \emph{portanto} introduz conclusão decorrente da informação anterior.
\end{questao}

\begin{questao}{\nivel{N3} 01.37 — FCC/referência}
Em ``A unidade revisou a norma antes de publicá-la'', o pronome \emph{la} retoma:
\begin{enumerate}[label=\Alph*)]
\item unidade;
\item norma;
\item revisão;
\item publicação;
\item sujeito oculto.
\end{enumerate}
\end{questao}

\begin{questao}{\nivel{N3} 01.38 — FGV/formalidade}
Assinale a reescrita mais adequada a relatório técnico formal.
\begin{enumerate}[label=\Alph*)]
\item A galera do setor vacilou feio nos controles.
\item O setor deu uma mancada na conferência.
\item Foram identificadas deficiências nos controles de conferência do setor.
\item O pessoal não fez direito o trabalho.
\item A coisa ficou meio errada no setor.
\end{enumerate}
\end{questao}

\begin{questao}{\nivel{N3} 01.39 — Cebraspe/Certo ou Errado}
A substituição de ``pode causar'' por ``causará'' aumenta o grau de certeza da afirmação e, portanto, pode alterar o sentido mesmo que a nova frase permaneça gramatical.
\end{questao}

\begin{questao}{\nivel{N3} 01.40 — FCC/redução de oração}
A reescrita ``Ao concluir a análise, a equipe emitiu parecer'' pressupõe, na leitura mais natural, que:
\begin{enumerate}[label=\Alph*)]
\item sujeitos das duas ações são distintos;
\item a equipe concluiu a análise e emitiu o parecer;
\item o parecer concluiu a análise;
\item não existe relação temporal;
\item a análise foi concluída após o parecer.
\end{enumerate}
\end{questao}

\subsection*{N4 — Auditoria / avançado}

\begin{questao}{\nivel{N4} 01.41 — Cebraspe/caso integrado}
Original: ``Os contratos com indícios de sobrepreço serão prioritariamente auditados.'' Julgue o item:

A reescrita ``Somente contratos com indícios de sobrepreço serão auditados'' preserva integralmente o sentido, pois \emph{prioritariamente} e \emph{somente} restringem o mesmo conjunto.
\end{questao}

\begin{questao}{\nivel{N4} 01.42 — FGV/reescrita complexa}
Original: ``Embora o controle estivesse formalmente instituído, não havia evidência de sua execução periódica.'' A melhor reescrita é:
\begin{enumerate}[label=\Alph*)]
\item Como o controle estava instituído, havia evidência periódica.
\item Apesar de formalmente instituído, o controle não tinha sua execução periódica comprovada.
\item O controle não estava instituído nem executado.
\item A evidência periódica impediu a instituição do controle.
\item Se o controle fosse instituído, haveria evidência.
\end{enumerate}
\end{questao}

\begin{questao}{\nivel{N4} 01.43 — FCC/escopo da negação}
A frase ``A auditoria não identificou todos os pagamentos irregulares'' permite concluir corretamente que:
\begin{enumerate}[label=\Alph*)]
\item nenhum pagamento irregular foi identificado;
\item pelo menos um pagamento irregular não foi identificado, sem excluir que outros tenham sido;
\item todos os pagamentos foram regulares;
\item a auditoria identificou exatamente um irregular;
\item não existiam pagamentos irregulares.
\end{enumerate}
\end{questao}

\begin{questao}{\nivel{N4} 01.44 — Cebraspe/Certo ou Errado}
Em texto de auditoria, substituir ``há risco de ocorrência'' por ``a ocorrência é certa'' não é mera simplificação estilística, pois transforma possibilidade em certeza.
\end{questao}

\begin{questao}{\nivel{N4} 01.45 — FGV/coesão e lógica}
Original: ``Os logs estavam incompletos; por isso, não foi possível confirmar integralmente o cumprimento do SLA.'' Qual alternativa inverte a relação lógica?
\begin{enumerate}[label=\Alph*)]
\item Como os logs estavam incompletos, não foi possível confirmar integralmente o SLA.
\item A incompletude dos logs impediu a confirmação integral do SLA.
\item Não foi possível confirmar integralmente o SLA, pois os logs estavam incompletos.
\item Os logs estavam incompletos, embora tenha sido impossível confirmar integralmente o SLA.
\item A confirmação integral ficou prejudicada em razão da incompletude dos logs.
\end{enumerate}
\end{questao}

\begin{questao}{\nivel{N4} 01.46 — Cebraspe/Certo ou Errado}
A frase ``O auditor recomendou que se revisassem os acessos'' admite próclise em razão da estrutura subordinada introduzida por \emph{que}.
\end{questao}

\begin{questao}{\nivel{N4} 01.47 — FCC/clareza}
Qual redação é mais precisa para um achado?
\begin{enumerate}[label=\Alph*)]
\item Existem vários problemas sérios no setor.
\item O setor aparentemente não funciona muito bem.
\item Em 17 de 40 contas analisadas, não havia registro de revisão no período exigido pela política.
\item A segurança é fraca e precisa melhorar.
\item Tudo indica que faltam controles.
\end{enumerate}
\end{questao}

\begin{questao}{\nivel{N4} 01.48 — FGV/pontuação e ambiguidade}
A frase ``Não foram analisados, por falta de acesso, os registros do período anterior'' usa as vírgulas para:
\begin{enumerate}[label=\Alph*)]
\item separar sujeito do verbo;
\item isolar expressão intercalada de causa;
\item indicar oração adjetiva restritiva;
\item marcar vocativo;
\item eliminar a necessidade de regência.
\end{enumerate}
\end{questao}

\begin{questao}{\nivel{N4} 01.49 — Cebraspe/Certo ou Errado}
Em ``A equipe, ao revisar os dados, identificou duplicidades'', a retirada das vírgulas sem qualquer outra alteração prejudicaria a organização sintática do segmento intercalado.
\end{questao}

\begin{questao}{\nivel{N4} 01.50 — FGV/revisão técnica}
Assinale a versão mais adequada de ``Foi verificado pela equipe que tinha um monte de conta velha ainda ativa'' para relatório formal.
\begin{enumerate}[label=\Alph*)]
\item A equipe viu que tinha conta velha demais.
\item Verificou-se muitas contas velhas ativas.
\item A equipe identificou diversas contas de usuários desligados que permaneciam ativas.
\item Foi visto que as contas estavam meio antigas.
\item As contas velhas, a equipe viu elas ativas.
\end{enumerate}
\end{questao}